\subsection{Проблемы извлечения, хранения и предоставления ончейн-данных}

При прямой работе прикладных систем с узлом Bitcoin Core возникают трудности,
связанные как с моделью хранения данных в блокчейне, так и с эксплуатационными
ограничениями RPC-интерфейса. Узел предоставляет корректные первичные данные,
однако не предназначен для эффективного обслуживания большого количества
сложных прикладных запросов, например выборок транзакций по адресам,
исторических балансов или актуальных наборов \gls{utxo}.

Основные проблемы извлечения и хранения ончейн-данных можно свести к
следующим положениям:
\begin{enumerate}
    \item данные представлены в нормализованном для протокола, но не для
    прикладного использования виде;
    \item вычисление производных сущностей, таких как адресные балансы,
    требует последовательной обработки значительных объемов истории;
    \item повторный пересчет агрегатов при каждом пользовательском запросе
    приводит к неприемлемым затратам времени и ресурсов;
    \item reorg требует корректного отката или перерасчета части ранее
    сохраненных представлений;
    \item mempool создает отдельный класс временных данных, которые должны
    наблюдаться и выдаваться обособленно от подтвержденной цепи.
\end{enumerate}

С точки зрения хранения значимым является вопрос согласованности. Если блок,
его транзакции, набор текущих \gls{utxo} и адресные балансы обновляются
несогласованно, то прикладная система теряет достоверность результатов. По этой
причине система индексации должна опираться на транзакционную модель записи,
идемпотентную обработку повторных событий и фиксированные правила перехода
между состояниями данных.

Проблема предоставления данных также требует отдельного решения. Для внешних
клиентов необходим интерфейс, который скрывает детали внутреннего формата
блокчейна и предоставляет стабильные прикладные операции: получение блоков,
транзакций, балансов адресов, mempool-транзакций и статуса задач индексации.
Следовательно, эффективная работа с ончейн-данными предполагает наличие
специализированного промежуточного слоя между блокчейн-узлом и прикладными
потребителями.
